%% ma: L12-B11-0019
%% lop: 12
%% chuong: 4
%% bai: 11
%% dang: 12.11.14
%% loai: TN4
%% muc: TH
%% dap_an: "D"
%% nguon: "Ảnh giáo viên gửi 24/10/2026 (M06): Kỳ thi thử đợt IV, Thi thử TN THPT 2027 (Biên soạn: Aimbet), Phần 1 Câu 1"
%% kien_thuc: "Bài 11 (nguyên hàm của 1/x, nguyên hàm thỏa điều kiện cho trước)"
%% ghi_chu: "Đề gốc không có phương án đúng (F(2) = ln 6 + 2, phương án D gốc là ln 5 + 2 = F(1)); đã sửa phương án D thành ln 6 + 2"
%% trang_thai: chinh-thuc
\begin{ex}
Cho $F(x)$ là một nguyên hàm của hàm số $f(x)=\dfrac{1}{x+4}$ thỏa mãn $F(-3)=2$. Giá trị của $F(2)$ bằng
\choice{$5\ln2$}{$2\ln5$}{$\ln2+5$}{\True $\ln6+2$}
\loigiai{$F(x)=\ln|x+4|+C$. $F(-3)=\ln1+C=2\Rightarrow C=2$. Vậy $F(2)=\ln6+2$.}
\end{ex}
